\chapter{Introduction}
Radio astronomy enables the exploration of celestial objects by their emission in radio frequency, revealing phenomena like jets from active galactic nuclei, halos and relics in galaxy clusters, and magnetic fields in the interstellar medium \cite{Chibueze2023}. Radio-wavelength observations probe physical processes like synchrotron emission and the 21\,cm line of neutral hydrogen. Over the past decades, radio astronomy has expanded exponentially due to technological advances in digital signal processing and new scientific drivers. Large interferometric arrays (e.g., ALMA, the Jansky VLA, the MeerKAT telescope, and the SKA) achieve high angular resolution and sensitivity by combining large effective apertures \cite{Chibueze2023}. For example, the 64-dish MeerKAT array in South Africa, which is a precursor to the SKA, performs crucial surveys in continuum imaging, deep HI mapping, and pulsar timing \cite{Jonas2018}. All these science programs depend importantly on the correlator being able to mix all antennas' signals. Wang et al. (2024) observe that a correlator is "one of the most important data handling components of a radio telescope array," pointing to its central role \cite{Wang2024}. Wang et al., for instance, describe a 192-input ROACH2+GPU correlator employed in the Tianlai 96-antenna pathfinder array, illustrating contemporary correlator scalability \cite{Wang2024}. The correlator's products (visibilities) are the data products on which imaging and analysis of the radio sky become possible.

Interferometry combines the signals from multiple antennas to simulate a very large aperture. Each antenna pair (baseline) measures a complex visibility, which is the Fourier component of the sky brightness sampled by that baseline. In an \emph{FX} correlator, for instance, each antenna’s time-stream voltage $E_i(t)$ is first digitised and Fourier-transformed to produce a spectrum $X_i(f)$ (the F-engine operation). The correlator then computes the product of one antenna’s spectrum with the complex conjugate of another’s,

\[
  V_{ij}(f) \;=\; X_i(f)\,X_j^*(f),
\]

yielding the complex visibility $V_{ij}(f)$ for antennas $i$ and $j$ at frequency $f$ \cite{Wang2024}. Equivalently, $V_{ij}(f)$ is the Fourier transform of the time-averaged cross-correlation $R_{ij}(\tau)=\langle E_i(t)E_j^*(t+\tau)\rangle$. The set of visibilities across many baselines samples the Fourier transform of the sky brightness (via the van Cittert–Zernike theorem), enabling image reconstruction by inverse Fourier techniques. (By contrast, an \emph{XF} correlator would form the time-domain cross-correlation $R_{ij}(\tau)$ first and then apply a Fourier transform to obtain $V_{ij}(f)$.) 

A notable example is the MeerKAT telescope’s Correlator/Beamformer (CBF) subsystem, which implements an FX architecture using CASPER FPGA hardware \cite{Jonas2018}. Signals from the 64 dishes (dual-polarised) are first channelised by digital polyphase filterbanks. After channelisation, a “corner-turn” is performed and the data are sent via 10-Gb Ethernet to the X-engines. The X-engines then carry out all cross-multiplications in real time on FPGA boards \cite{Jonas2018}, computing the thousands of visibilities needed for imaging. As described by Jonas (2018), MeerKAT’s correlator handles on the order of thousands of baselines per polarisation, processed continuously to build up images \cite{Jonas2018}.

    \section{Correlator Architectures: XF vs FX}
    Digital correlators can be categorised as \emph{XF} or \emph{FX} based on the order of operations \cite{Wang2024}. In an \textbf{XF correlator}, the time-domain signals from each antenna pair are aligned (delay-tracked) and multiplied for each time lag, producing a lag-spectrum which is then Fourier-transformed into the frequency domain. XF designs can naturally accommodate very fine spectral resolution and handle large delay ranges, but the number of multipliers required grows rapidly with the number of antennas (baselines). In an \textbf{FX correlator}, each antenna’s time-stream is first Fourier-transformed (typically using a polyphase filterbank) into frequency channels, and then the frequency-domain data from antenna pairs are multiplied channel by channel. FX correlators exploit efficient FFT algorithms (cost $\sim N_{\rm chan}\log N_{\rm chan}$ per antenna) and are well suited to parallel hardware (FPGAs or GPUs) \cite{Wang2024,Jonas2018}. Modern arrays like MeerKAT and planned SKA-mid favor the FX approach for its scalability \cite{Wang2024}. 
    
    Summary of the key differences between the XF and FX approaches:
    
    \begin{itemize}
        \item \textbf{XF correlator}: Cross-multiplication in time domain first, Fourier transform second. Good for extremely high spectral resolution; correlator hardware must handle long time buffers for each baseline.
     
        \item \textbf{FX correlator}: FFT (F-engine) on each antenna’s stream first, then cross-multiply in frequency domain (X-engine). Efficient for large numbers of antennas, easier to distribute among processors, and matches well to modern FPGA/GPU pipelines \cite{Wang2024}.
    \end{itemize}
    
    \section{Signal Chain and Processing}
    A radio telescope’s \emph{signal chain} refers to the entire processing path from the antenna feed to the final data output. After the dish collects the electromagnetic waves, the first stage is a low-noise amplifier (LNA) to boost the signal and reduce system noise. The signal is then down-converted from radio frequencies (RF) to an intermediate frequency (IF) or directly to baseband using mixers and local oscillators. Careful RF filtering is applied to suppress out-of-band interference. In the digital era, the down-converted signal is fed to an analog-to-digital converter (ADC). For example, the CASPER ADC16x250-8 board digitises 16 analog inputs at 8 bits and 250 Msps, covering 125 MHz of bandwidth per channel \cite{Wang2024}. All ADCs are phase-locked to a common 10\,MHz reference and pulse-per-second (PPS) signal (from a GPS-disciplined clock) to ensure time-alignment across antennas. In our design, the four dual-polarised antenna outputs (8 analog streams) will be handled by one or two such ADC boards, each configured for 8 or 16 inputs. 
    
    The moment the data are digitised, they come into the correlator's \emph{F-engine}. The F-engine channels the wideband data using a digital polyphase filter bank (PFB). A PFB employs a finite impulse response (FIR) window before an FFT, producing many narrow frequency channels of flat passbands and high out-of-band attenuation. This suppresses spectral leakage and improves imaging quality. For example, \cite{Wang2024} illustrate a design that consists of a 1024-tap PFB followed by a 1024-point FFT \cite{Wang2024}. Under this case, the 36-bit (18+18) complex PFB outputs were scaled and quantised to 4 bits for real and imaginary parts in order to reduce data volume \cite{Wang2024}. After channelisation, each frequency sub-band stream is time-stamped and ready for the X-engine.
    
    Before cross-multiplication, a corner-turn (data transpose) is performed to regroup data by frequency rather than by antenna/time. In practice, this involves buffering the channelised data in FPGA memory, then re-reading it so that each packet contains all antenna data for one or more frequency channels. These packets (carrying a few channels for one antenna) are sent via high-speed Ethernet (typically 10\,GbE) with headers indicating antenna ID, channel ID, time stamp, etc. The switch/router then distributes packets to the X-engine compute nodes according to the channel. 
    
    In the \emph{X-engine}, the actual correlation calculations occur. Each GPU or CPU node receives UDP packets containing spectra from all antennas for its assigned channels. The node then performs complex multiply-and-accumulate operations across all antenna pairs for those channels. For example, the xGPU library runs thousands of CUDA threads per GPU to compute all baselines in parallel. The output of the X-engine is the visibility data: complex values $V_{ij}(f_k)$ for every antenna pair $(i,j)$, polarisation combination, and frequency channel $f_k$. These visibilities are typically integrated over a short time interval and then recorded to disk or passed to an imaging pipeline. Note that autocorrelations ($i=j$) produce the power spectrum of each antenna, useful for system monitoring and total-power measurements. All together, for $N$ antennas with two polarisations each, the correlator computes on the order of $(2N)(2N-1)/2$ complex products per channel per integration.
    
    \section{Radio Telescope Orientation Software (RTOS)}
    An interferometric array also requires precise control of the antennas. Modern observatories use a unified control system (often called CAM: Control and Monitoring). For example, the MeerKAT array’s CAM is built on the Karoo Array Telescope Communication Protocol (KATCP) \cite{Schwartz2023}. KATCP allows software to configure hardware modules (antenna drives, receivers, etc.) through a text-based protocol. In the MeerKAT/AVN framework, the CAM includes modules such as the Station Controller (overall scheduler), Antenna Steering Control, Receiver Control, and Environmental Monitors \cite{Schwartz2023}. These modules handle tasks like commanding the dish azimuth/elevation drives, tuning the receiver’s local oscillator, switching filters, and monitoring temperatures or voltages. The RTOS also orchestrates calibration signals: for instance, periodic noise-diode injections into the receiver can be synchronised by the software to solve for system gains. 
    
    The Antenna Control Subsystem (ACS) calculates the real-time pointing commands. It takes source coordinates (RA/Dec) or ephemeris data and computes the required mount angles, accounting for Earth rotation, refraction, precession, and flexure. Encoder readings from the servos provide feedback to ensure accurate tracking. In practice, the ACS may apply a pointing model (with terms for alignment errors and tilt) derived from calibration observations. Moreover, the RTOS must distribute timing signals: all antennas’ ADC samplers must share the same clock and PPS. In VLBI, for example, each telescope stamps its data with a GPS or maser reference and the signals are correlated later by matching the timestamps \cite{Schwartz2023}. In a real-time array like Nooitgedacht, a common clock is shared, but the RTOS still issues a PPS to mark integration boundaries so that all data can be coherently combined.
    
    \section{The Nooitgedacht Telescope Farm and Project Significance}
    The North-West University Nooitgedacht facility now hosts four 3.7-meter radio dishes arranged in a compact configuration \cite{Chibueze2023}. The site (approximately 35\,km from Potchefstroom) is in a rural area with relatively low radio frequency interference \cite{Chibueze2023}. Each dish has a wide primary beam (of order a few degrees at L-band), which allows mapping of large sky regions. At 1–18\,GHz, the array will resolve sources on the order of several arcminutes (given baselines up to $\sim$50\,m). This is sufficient to image strong southern radio sources and extended structures (e.g.\ supernova remnants, radio lobes of nearby AGN). The array’s large field of view also makes it useful for surveying or monitoring variable sources and transients.
    
    However, when the dishes were installed around 2022, the array lacked a digital correlator. Without a correlator, the antennas could only operate as independent dishes. Implementing a real correlator will enable true interferometric operation. The new correlator must digitise 125\,MHz (or more) of bandwidth from each dish and compute all pairwise visibilities in both polarisations. In our case, 4 dual-polarised antennas mean 8 data streams, yielding $8\times7/2=28$ cross-correlation products plus 8 auto-correlations per frequency channel. The correlator design thus must handle at least these 36 complex products at the target sample rate.
    
    Constructing the correlator brings multiple benefits. Scientifically, it allows the array to form high-resolution synthesised images and spectra. Even with only four dishes, one can form useful interferometric beams and measure visibilities for strong sources. For example, dedicated monitoring of bright pulsars could track their flux and polarisation changes, and large-scale Galactic surveys (e.g.\ mapping HII regions or diffuse emission) become feasible. The correlator also opens the possibility of Very Long Baseline Interferometry (VLBI) by recording time-stamped data for later correlation with other telescopes. Technologically, this project provides a valuable training platform. Students will gain hands-on experience in FPGA programming, signal chain design, and high-performance computing (GPUs). 
    
    The significance extends to South African radio astronomy infrastructure. South Africa has invested heavily in MeerKAT and SKA, and building local expertise is a strategic goal. SARAO notes that the current MeerKAT correlator is implemented on FPGA SKARAB boards, while a new GPU-based correlator is under development for the MeerKAT extension \cite{SARAO2024}. This reflects a trend toward hybrid FPGA/GPU solutions for flexibility and cost-effectiveness. The Nooitgedacht correlator will serve as an independent demonstration of such technology in a South African context, and it complements large observatories by providing a nearby array for system tests and student projects.
    
    \section{Motivation and Objectives}
    The primary objective of this thesis is to design, implement, and validate a digital correlator for the Nooitgedacht radio telescope array. Specifically, we will: select and integrate the analog front-end and ADC hardware; develop FPGA firmware for the F-engine (channelizer and packetiser); establish the 10-GbE network infrastructure; develop the GPU or CPU code for the X-engine; and integrate the correlator with the existing RTOS. The correlator must meet the scientific requirements of the array (processing 1–18\,GHz band, 125\,MHz instantaneous bandwidth, dual polarisation) and deliver stable visibilities. We will also characterise the performance (sensitivity, stability) of the system with test signals and on-sky observations.
    
    \begin{itemize}
        \item \textbf{Enabling Interferometric Imaging:} Endow the Nooitgedacht array with the ability to perform full aperture-synthesis imaging and spectroscopy, exploiting the scientific potential of the array.

        \item \textbf{Training and Education:} Provide a hands-on setting for researchers and students to acquire digital signal processing, FPGA/GPU programming, and radio interferometry expertise.
    
        \item \textbf{Technology Demonstration:} Explore the utilization of FPGA (CASPER) and GPU hardware within a small-sized correlator, which will be utilized for giving direction towards large array designs.

        \item \textbf{Augment Infrastructure:} Deliver beneficial capability to South Africa's radio astronomy infrastructure, complementing large facilities (MeerKAT/SKA) and providing a local testbed for observation and verification.
    \end{itemize}
    
    To achieve these goals, the design emphasises modularity and scalability. Wherever possible, we will use open-source FPGA designs (e.g.\ CASPER DSP blocks) and commodity GPU hardware. The correlator firmware and software will be written with flexibility in mind, allowing future upgrades (e.g.\ adding a beamforming mode or implementing fringe-rotation). The correlator will be integrated into the observatory’s control system so that observations can be automated and data can be archived in standard formats. By the end of this work, the Nooitgedacht array will operate as a complete interferometer, providing real scientific data and training experience.
    
    \section{Thesis Outline}
    Chapter 2 provides a comprehensive review of correlator architectures and digital backend design from the recent literature. Topics include FX vs XF designs, polyphase filterbanks, packetised data networks, and example implementations in existing telescopes. This chapter also describes the system design for the Nooitgedacht correlator: the analog front-end and receiver chain, the choice and configuration of ADC modules, and the FPGA-based F-engine firmware. The correlator backend (X-engine): it details the network setup, the GPU implementation of the cross-correlation, and the data formatting for storage is also covered, along with the system integration and testing: the methods used to validate the correlator (injection of test signals, on-sky calibrators) and the performance metrics achieved.
    
    In summary, this introduction has established the context and key concepts for the Nooitgedacht correlator project. By outlining the scientific motivation and reviewing relevant technology, it lays the groundwork for the detailed design and implementation that follow. The next chapters will build on this foundation to realise a fully functional interferometer, thereby enabling South African researchers to fully exploit the Nooitgedacht array.